\documentclass{ohm_project_description}
\usepackage[english]{babel}

\projecttitle{Development of a Mobile Mini-Robot}
\projectauthor{Mobile Robotics Lab}
\projectdate{2026}

\begin{document}

\maketitle
\thispagestyle{fancy}

In this project, a mobile mini-robot will be designed, built, and tested as a
development and training platform. Starting from an existing robot concept, an
independent, mechanically and electrically robust variant will be created that can be
reproduced cost-effectively – with material costs under €1,000 per unit. The robot
features an omnidirectional drive with four Mecanum wheels, a rechargeable power
supply with multiple on-board voltages, and a powerful single-board controller with
active cooling. A 360° LiDAR is integrated for environment sensing; a depth camera
can optionally be added. The robot is operated either via Bluetooth controller or a
browser-based web interface on a tablet.

\section*{Work Packages}
\begin{itemize}[leftmargin=0.5cm]
    \setlength\itemsep{.1em}
    \item Design of the mechanical structure (multi-deck aluminum frame, component
          mounting with low center of gravity)
    \item Selection, procurement, and electrical integration of components (drive,
          power supply with DC-DC conversion, charging concept)
    \item Commissioning of the motor controller and omnidirectional drive (Mecanum
          wheels, modular driver boards)
    \item Integration and testing of sensors and output devices (360° LiDAR, display,
          speaker, optional depth camera)
    \item Implementation of user interfaces (Bluetooth controller, web interface, data
          recording to external storage)
    \item Documentation for reproduction and cost analysis for later series production
\end{itemize}

\section*{Requirements}
\begin{itemize}[leftmargin=0.5cm]
    \setlength\itemsep{.1em}
    \item Interest in mobile robotics and hands-on system integration
    \item Basic knowledge of electrical engineering (circuit design, power supply)
    \item Basic knowledge of mechanical engineering (structural design, CAD
          advantageous)
    \item Basic programming skills (Python or C++ helpful)
    \item Ability to work in a team and enjoy interdisciplinary collaboration
    \item Suitable for students of BEEI, BEME, and related degree programmes
\end{itemize}

\vspace{0.5cm}
This topic can be completed as a \textbf{project work} (4 students, approx.\ 4 months)
subject to agreement.

\vfill
\textcolor{ohm_red}{\rule{\linewidth}{0.4mm}}
\textbf{\textcolor{ohm_red}{Mobile Robotics Lab}} \\
\begin{tabular}{@{}ll}
\textbf{Supervisor:} & Prof.\ Dr.\ Christian Pfitzner \\
\textbf{E-Mail:}     & \href{mailto:christian.pfitzner@th-nuernberg.de}{christian.pfitzner@th-nuernberg.de} \\
\end{tabular}

\end{document}
